\documentclass[11pt]{article}
\usepackage{fontspec}
\usepackage[spanish]{babel}
%\usepackage[utf8]{inputenc}
\usepackage{listings}
\usepackage{graphicx}
\graphicspath{{../Imagenes/}}

\usepackage{titlepic}

%%% Tablas
\usepackage{longtable}
\usepackage{tabularx}
\usepackage{float}
\usepackage{adjustbox}
\usepackage{booktabs}
\usepackage{multirow}
\renewcommand{\arraystretch}{1.7}

%%%%%%%%%% Contador HU
\newcounter{HUCounter}
\newcommand{\hu}[1]{\refstepcounter{HUCounter}\textbf{\rmfamily HU-\theHUCounter}\label{#1}}

\begin{document}

\begin{titlepage}
\centering
\vspace{4.5cm}
{\scshape\LARGE Historias de Usuario \par}
\vspace{1.5cm}

\includegraphics[width=10cm]{Logo}

\vspace{3cm}
{\scshape\large \par}
\vspace{1cm}

{Miguel Albertí Pons\\
Sofía Almeida Bruno\\
Pedro Manuel Flores Crespo\\
María Victoria Granados Pozo\\
Lidia Martín Chica
\par}

\end{titlepage}
\newpage

\section{Listado de Historias de Usuario}
En el apartado de prioridad hemos escogido una escala del 1 al 6 siendo el 1: máxima prioridad y el 6: prioridad mínima.
  \begin{longtable}{p{0.12\linewidth}p{0.7\linewidth}p{0.15\linewidth}p{0.15\linewidth}}
    \toprule
    \textbf{Id HU} & \textbf{Título} & \textbf{Estimación} & \textbf{Prioridad}\\
    \midrule
     \hu{} & Un usuario quiere registrarse en la aplicación para poder planificar su viaje & 2 & 1\\
     \hu{} & Un usuario quiere identificarse en la aplicación para  acceder a su contenido & 1 & 1\\
     \hu{} & Un usuario quiere ver las publicaciones para saber cuáles le pueden interesar & 2 & 2\\
    \hu{} & Un usuario quiere ofertar un viaje/actividad para compartir su experiencia & 3 & 2\\
    \hu{} & Un usuario quiere valorar un viaje que ha realizado para mostrar su opinión & 2 & 4\\
    \hu{} & Un usuario quiere comentar en una publicación para preguntar una duda o aportar información & 2 & 4\\ 
    \hu{} & Un usuario quiere ver los comentarios de una publicación para conocer las experiencias y opiniones de otros usuarios & 1 & 4\\ 
    \hu{} & Un usuario quiere ver las valoraciones de una publicación para tenerlo en cuenta a la hora de tomar una decisión & 1 & 4\\ 
    \hu{hu:filtrar} & Un usuario quiere filtrar la búsqueda de las publicaciones para encontrar una que se adapte a sus intereses & - & -\\ 
    \hu{hu:personalizar} & Un usuario quiere personalizar un paquete para adaptarlo a sus gustos & - & - \\ 
    \hu{} & Un usuario quiere seleccionar un viaje/actividad para realizarlo & 3/2 & 2\\ 
    \hu{} & Un usuario quiere guardar un viaje/actividad para verlo en otro momento & 1 & 5\\ 
    \hu{}. & Un usuario quiere ver su historial de viajes guardados/lista de deseos para elegir entre ellos & 1 & 5\\
    \hu{hu:maleta} & Un usuario quiere gestionar su maleta & - & -\\
    \hu{} & Un usuario quiere grabar un audioguía sobre su ciudad para compartirla con otros usuarios & 2 & 6\\
    \hu{} & Un usuario quiere escuchar un audioguía sobre la ruta que va a realizar para informarse sobre la ciudad & 3/2 & 6\\

    
    \bottomrule
\end{longtable}

\subsection{Historias que se han dividido en varias}

  \begin{longtable}{p{0.13\linewidth}p{0.67\linewidth}p{0.15\linewidth}p{0.15\linewidth}}
    \toprule
    \textbf{Id HU} & \textbf{Título} & \textbf{Estimación} & \textbf{Prioridad}\\
    \midrule
    
    \textbf{HU-\ref{hu:filtrar}} & Un usuario quiere filtrar la búsqueda de las publicaciones para encontrar una que se adapte a sus intereses & - & -\\ 
	\textbf{HU-\ref{hu:filtrar}.1} & Un usuario quiere filtrar la búsqueda de las publicaciones por lugar & 1/2 & 3\\ 
	\textbf{HU-\ref{hu:filtrar}.2} & Un usuario quiere filtrar la búsqueda de las publicaciones por fecha & 1/2 & 3\\ 
	\textbf{HU-\ref{hu:filtrar}.3} & Un usuario quiere filtrar la búsqueda de las publicaciones por tipo de publicación (paquete, ruta, actividad, evento) & 1/2 & 3\\ 
	\textbf{HU-\ref{hu:filtrar}.4} & Un usuario quiere filtrar la búsqueda de las publicaciones por tipo de actividad (cultural, deportiva, etc)) & 1/2 & 3\\ 
	\textbf{HU-\ref{hu:filtrar}.5} & Un usuario quiere filtrar la búsqueda de las publicaciones para que aparezcan las más visitadas & 1/2 & 4\\ 
	\textbf{HU-\ref{hu:filtrar}.6} & Un usuario quiere filtrar la búsqueda de las publicaciones para que aparezcan las más realizadas & 1/2 & 4\\ 
	\midrule
	\textbf{HU-\ref{hu:personalizar}} & Un usuario quiere personalizar un paquete para adaptarlo a sus gustos y presupuesto & - & - \\
	\textbf{HU-\ref{hu:personalizar}.1} & Un usuario quiere eliminar una actividad de un paquete & 1/2 & 3 \\ 
	\textbf{HU-\ref{hu:personalizar}.2} & Un usuario quiere añadir una actividad de un paquete & 1 & 3 \\ 
	\textbf{HU-\ref{hu:personalizar}.3} & Un usuario quiere modificar el orden de las actividades de un paquete & 3/2 & 4 \\
	\midrule 
	\textbf{HU-\ref{hu:maleta}} & Un usuario quiere gestionar su maleta & - & -\\
	\textbf{HU-\ref{hu:maleta}.1} & Un usuario quiere obtener una recomendación del vestuario que debe llevar según el pronóstico meteorológico & 5 & 6\\
	\textbf{HU-\ref{hu:maleta}.2} & Un usuario quiere obtener una recomendación del vestuario que debe llevar según el tipo de actividad que va a realizar & 4 & 6\\
	\textbf{HU-\ref{hu:maleta}.3} & Un usuario quiere obtener una recomendación del vestuario que debe llevar según la duración del viaje & 4 & 6\\
    \bottomrule
  \end{longtable}

\section{Cálculo de la velocidad del equipo}
Partimos de	un equipo de desarrollo	
  formado	
  por	
  5	
  programadores	
  que	
  van	
  a	
  dedicar	
  un	
  60\%	
  de	
  su	
  trabajo	
  al	
  proyecto, ya que el resto del tiempo lo utilizan para avanzar otros proyectos.	
  

La	
  duración	
  de	
  cada	
  una	
  de	
  las	
  iteraciones	
  que	
  vamos	
  a	
  realizar	
  en	
  el	
  proyecto	
  van	
  a	
  ser	
  de	
  2	
  semanas.	
  
La	
  estimación	
  realizada	
  del	
  esfuerzo	
  de	
  cada	
  una	
  de	
  las	
  historias	
  de	
  usuario	
  se	
  ha	
  expresado	
  en	
  días	
  ideales	
  de	
  programación.	
  En	
  nuestro	
  
entorno	
  de	
  trabajo	
  estimamos	
  que	
  un	
  día	
  ideal	
  de	
  programación	
  se	
  va	
  a	
  corresponder	
  con	
  2	
  a	
  3	
  días	
  reales	
  de	
  trabajo.	
  

La	
   duración	
  de	
  una	
  iteración	
  va	
  a	
  ser:	
  

1	
iteración	
  =	
  2	
  semanas	
  =	
  6 días	
  reales	
  

La	
   velocidad	
  del	
  equipo	
  de	
  desarrollo	
  medido	
  en	
  punto	
  de	
  historia	
  es:	
  	
  

5	
  programadores	
  *	
  6	
  = 30	
  	
  días	
  reales	
  por	
  iteración $\Longrightarrow$	
  de	
  10	
  a	
  15	
  PH	
  por	
  iteración.	
  	
  	
  	
  	
  
Se	
  ha	
  decidido	
  usar	
   12	
  Puntos	
  de	
  historia	
  como	
  la	
  velocidad	
  estimada	
  del	
  equipo.	
  

\section{Descripción de las entregas}
Esfuerzo total del proyecto = 40 PH \\
Velocidad del equipo = 12 PH (por iteración) 

\end{document}
